\section{The Exponential Series}

Consider the function $f \colon \CC \to \CC$
defined by 
\begin{equation}\label{eq:def_exp}
f(z) = \sum_{k=0}^{+\infty} \frac{z^k}{k!}.
\end{equation}
The function is defined on the entire $\CC$ since (thanks to the ratio test)
it is easy to verify that the radius of convergence of this series is
$R=+\infty$.

Our goal will be to prove that $f(z)=e^z$ for every $z\in \CC$.
In fact, one could take~\ref{eq:def_exp} as the definition 
of the complex exponential and define the real exponential 
function and the trigonometric functions accordingly.

\begin{theorem}[Properties of the Exponential Series]
\label{th:exp_complesso}%
Let $f$ be the function defined by ~\eqref{eq:def_exp}.
We have:
\begin{enumerate}
\item
$\displaystyle f(0) = 1$;

\item
$\displaystyle f(\bar z) = \overline{f(z)}$;

\item
for every $z,w \in \CC$
\[
  f(z+w) = f(z) \cdot f(w);
\]

\item
for every $z\in \CC$, $f(z) \neq 0$ and
\[
 f(-z) = \frac{1}{f(z)};
\]

\item the function $f\colon \CC \to \CC$ is continuous;

\item we have
\begin{equation}\label{eq:limite_exp_complesso}
   \lim_{z\to 0}\frac{f(z)-1}{z} = 1.
\end{equation}
\end{enumerate}
\end{theorem}
%
\begin{proof}
\mymark{*}
\begin{enumerate}
\item
The property $f(0)=1$ is obtained by direct verification (recall that $0^0=1$
and $0!=1$).

\item The property $f(\bar z) = \overline{f(z)}$ is obtained
by taking the limit of the following equality between partial sums
which exploits the properties of the conjugate of a sum and a product:
\[
\sum_{k=0}^{n} \frac{\bar z^k}{k!} = \overline{\sum_{k=0}^n \frac{z^k}{k!}}.
\]

\item
Consider the infinite matrix
\[
m_{k,j}  = \frac{z^k}{k!} \cdot \frac{w^j}{j!}.
\]
Then on one hand
\begin{align*}
 f(z+w)
 &= \sum_{n=0}^{+\infty} \frac{(z+w)^n}{n!}\\
  &= \sum_{n=0}^{+\infty} \frac{1}{n!}\sum_{k=0}^n \frac{n!}{k!(n-k)!} z^k\cdot
 w^{n-k}\\
 &= \sum_{n=0}^{+\infty} \sum_{k=0}^n m_{k, n-k}
\end{align*}
and on the other
\begin{align*}
 f(z) \cdot f(w)
 &= \sum_{k=0}^{+\infty}\frac{z^k}{k!} \sum_{j=0}^{+\infty}\frac{w^j}{j!}
 = \sum_{k=0}^{+\infty}\sum_{j=0}^{+\infty}\frac{z^k}{k!} \frac{w^j}{j!} \\
 &= \sum_{k=0}^{+\infty}\sum_{j=0}^{+\infty} m_{k,j}.
\end{align*}

The result follows from theorem~\ref{th:somma_Cauchy}.

\item
Since
\[
  1 = f(0) = f(z-z) = f(z) \cdot f(-z)
\]
we deduce that $f(z)\neq 0$ and $f(-z) =  1 / f(z)$.

\item
Continuity follows from the general result on the continuity of
the sum of a power series: theorem~\ref{th:continuita_somma_serie}.

\item
Directly from definition~\eqref{eq:def_exp}
we obtain
\[
 f(z) - 1 = \sum_{k=1}^{+\infty} \frac{z^k}{k!}
 = z \cdot \sum_{k=0}^{+\infty} \frac{z^k}{(k+1)!}.
\]
Now observe that the series $\sum \frac{z^k}{(k+1)!}$ has an infinite radius of
convergence and hence is absolutely convergent for every $z\in \CC$.
In particular, the sum of this series is continuous and
therefore we have
\[
\lim_{z\to 0}\frac{\exp(z) - 1}{z}
   =  \lim_{z\to 0} \sum_{k=0}^{+\infty} \frac{z^k}{(k+1)!}
   = \sum_{k=0}^{+\infty}\frac{0^k}{(k+1)!} = 1.
\]
\end{enumerate}
\end{proof}

\begin{theorem}[Coincidence Between Exponential Function and Exponential Series]%
\label{th:exp=ex}%
\label{th:serie_esponenziale}%
\index{function!exponential}%
\mymark{***}%
For every $x\in \RR$ we have
\[
  \sum_{k=0}^{+\infty} \frac{x^k}{k!} = e^x.
\]
\end{theorem}
%
\begin{proof}\mymark{**}
We want to apply theorem~\ref{th:isomorfismo}.
Let $f(x)=\sum_{k=0}^{+\infty} \frac{x^k}{k!}$, 
$f\colon \RR\to\RR$ we know,
from the previous theorem, that $f(x+y)=f(x)\cdot f(y)$.
We need to verify that $f$ is increasing.
First, note that if $x> 0$ the series defining
$f(x)$ has positive terms and the first term is $1$.
Thus for every $x> 0$ we have $f(x)>1$.
In particular, let $a=f(1)$, then $a>1$.
Moreover, if $y>x$ we have $f(y-x)>1$ and thus
\[
 f(y) = f(x+y-x) = f(x)\cdot f(y-x) > f(x).
\]
The function is therefore strictly increasing. 
Thus it is a monotonic homomorphism from the additive group $\RR$ to 
the multiplicative group $(0,+\infty)$.
Given $a=f(1)$, we know,
from what was seen in section~\ref{sec:esponenziale},
that there is a unique function with these properties and thus 
it must be: 
\[
  f(x)= a^x.
\]
We only need to prove that $a=e$.
But thanks to theorem~\ref{th:exp_complesso}
just proven, we know that as $x\to 0$ we have
\[
  \frac{f(x)-1}{x}\to 1.
\]
On the other hand, we know that $f(x)=a^x$ and thus
by the notable limit (corollary~\ref{cor:limite_notevole_e})
\[
\frac{f(x)-1}{x} = \frac{a^{x}-1}{x}
= \frac{e^{x \ln a}-1}{x \ln a} \cdot \ln a
\to \ln a.
\]
We deduce that $\ln a= 1$, which means $a=e$.
\end{proof}

With some effort, it would also be possible to prove that 
indeed $f(z)=e^z$ (in the previous theorem, we proved it 
only for $x\in \RR$). 
It would involve verifying that the function 
$\phi(x) = f\enclose{i\frac{x}{2\pi}}$
has the properties stated in theorem~\ref{th:omomorfismo_U}.
However, it will be simpler to use derivatives, so we will postpone this 
verification to the next chapter.

In corollary~\ref{cor:limite_notevole_ex}
we proved that 
for every $x\in \RR$ we have
\[
  e^x = \lim_{n\to +\infty} \enclose{1+\frac x n}^n.  
\]
The following theorem tells us that the exponential series $f(z)$ 
satisfies the same property for every $z \in \CC$.
Thus, it indeed provides an alternative proof 
of the identity $f(x) = e^x$ and a further indication 
of the coincidence between $f(z)$ and $e^z$.

\begin{theorem}[Notable Limit of Complex Exponential]%
\label{th:limite_notevole_esponenziale_complesso}%
\mymark{*}%
For every $z\in \CC$ it holds that
\[
  \sum_{k=0}^{+\infty} \frac{z^k}{k!} = \lim_{n \to +\infty} \enclose{1+\frac z n}^n  
\]
\end{theorem}
%
\begin{proof}
\mymark{*}
Using the binomial expansion, we observe that
\[
 \enclose{1+\frac z n}^n
 = \sum_{k=0}^n \binom{n}{k} \frac{z^k}{n^k}
 = \sum_{k=0}^n \frac{z^k}{k!} \cdot \frac{n!}{n^k\cdot (n-k)!}.
\]
For each $k\le n$, let
\begin{align*}
 c(n,k)
  &= \frac{n!}{n^k\cdot (n-k)!}
  = \frac{n \cdot (n-1) \cdot \ldots \cdot(n-k+1)}{n^k} \\
    &= \frac{n}{n}\cdot {\frac {n-1} n} \cdot \frac {n-2} {n} \cdot \ldots \cdot
  \frac{n-k+1}{n}
\end{align*}
we observe that $0\le c(n,k)\le 1$ as it is
a product of non-negative numbers less than or equal to $1$.
Moreover, for fixed $k$, we have $c(n,k)  \to 1$ as $n\to +\infty$
since each factor $\frac{n-j}{n}$ tends
to $1$ as $n\to +\infty$ (note that for fixed $k$, the number of factors $k$ is
fixed).

Let $z\in \CC$ be fixed and let
\[
  S = \sum_{k=0}^{+\infty} \frac{\abs{z}^k}{k!}.
\]
We know that the exponential series is absolutely convergent
for every $z\in \CC$ (since the radius of convergence is $+\infty$)
so $S$ is a real (finite) number.
Thus, by theorem \ref{th:coda} (the tail theorem), we know that for every
$\eps>0$
there exists $M$ such that
\[
   \sum_{k=M+1}^{+\infty} \frac{\abs{z}^k}{k!} < \eps.
\]

For fixed $k\le M$, since $c(n,k)\to 1$,
there exists $N_k > M$ such that for every $n>N_k$
we have $1-c(n,k) < \eps$
(recall that $c(n,k)\le 1$).
Then take
\[
  N=\max\ENCLOSE{N_k\colon k\le M}
\]
so that for every $n>N$ and for every $k\le M$, we have $0 \le 1-c(n,k) < \eps$.
Then, for every $n>N$, we can split the sum from $0$ to $n$ into two
sums from $0$ to $M$ and from $M+1$ to $n$:
\begin{align*}
\abs{\sum_{k=0}^n \frac{z^k}{k!} - \enclose{1+\frac z n}^n}
&= \abs{\sum_{k=0}^n \enclose{\frac{z^k}{k!} - c(n,k)\frac{z^k}{k!}}}
= \abs{\sum_{k=0}^n  (1-c(n,k))\frac{z^k}{k!}} \\
&\le \sum_{k=0}^n  (1-c(n,k))\frac{\abs{z}^k}{k!} \\
  &= \sum_{k=0}^{M} (1-c(n,k)) \frac{\abs{z}^k}{k!}
   + \sum_{k=M+1}^n (1-c(n,k)) \frac{\abs{z}^k}{k!} \\
&\le \sum_{k=0}^{M} \eps \frac{\abs{z}^k}{k!}
   + \sum_{k=M+1}^n \frac{\abs{z}^k}{k!} \\
&\le  \eps \sum_{k=0}^{+\infty} \frac{\abs{z}^k}{k!}
    + \eps \\
&\le \eps S + \eps
= \eps (S+1).
\end{align*}

Since $\eps>0$ was arbitrary, we have verified
through the definition that
\[
\sum_{k=0}^n \frac{z^k}{k!} - \enclose{1+\frac z n}^n \to 0
\]
that is
\[
\lim_{n\to +\infty} \enclose{1+\frac z n}^n = \sum_{k=0}^{+\infty} \frac{z^k}{k!}.
\]
\end{proof}

\begin{table}
\begin{center}
\begin{tabular}{r}
\ttfamily\footnotesize 2.7182818284 5904523536 0287471352 6624977572 4709369995 \\
\ttfamily\footnotesize   9574966967 6277240766 3035354759 4571382178 5251664274 \\
\ttfamily\footnotesize   2746639193 2003059921 8174135966 2904357290 0334295260 \\
\ttfamily\footnotesize   5956307381 3232862794 3490763233 8298807531 9525101901 \\
\ttfamily\footnotesize   1573834187 9307021540 8914993488 4167509244 7614606680 \\
\ttfamily\footnotesize   8226480016 8477411853 7423454424 3710753907 7744992069 \\
\ttfamily\footnotesize   5517027618 3860626133 1384583000 7520449338 2656029760 \\
\ttfamily\footnotesize   6737113200 7093287091 2744374704 7230696977 2093101416 \\
\ttfamily\footnotesize   9283681902 5515108657 4637721112 5238978442 5056953696 \\
\ttfamily\footnotesize   7707854499 6996794686 4454905987 9316368892 3009879312 \\
\ttfamily\footnotesize   7736178215 4249992295 7635148220 8269895193 6680331825 \\
\ttfamily\footnotesize   2886939849 6465105820 9392398294 8879332036 2509443117 \\
\ttfamily\footnotesize   3012381970 6841614039 7019837679 3206832823 7646480429 \\
\ttfamily\footnotesize   5311802328 7825098194 5581530175 6717361332 0698112509 \\
\ttfamily\footnotesize   9618188159 3041690351 5988885193 4580727386 6738589422 \\
\ttfamily\footnotesize   8792284998 9208680582 5749279610 4841984443 6346324496 \\
\ttfamily\footnotesize   8487560233 6248270419 7862320900 2160990235 3043699418 \\
\ttfamily\footnotesize   4914631409 3431738143 6405462531 5209618369 0888707016 \\
\ttfamily\footnotesize   7683964243 7814059271 4563549061 3031072085 1038375051 \\
\ttfamily\footnotesize   0115747704 1718986106 8739696552 1267154688 9570350354
\end{tabular}
\end{center}
\caption{The first 1000 decimal digits of the number $e$
calculated using the method employed in the proof
of theorem~\ref{th:approx_e}.
See the code on page~\pageref{code:compute_e}.}
\label{fig:cifre_e}
\index{$e$!decimal digits}
\index{digits!$e$}
\end{table}

\begin{theorem}[Approximation of the Exponential]
  \label{th:approx_exp}%
  \label{th:approx_e}%
  \index{$\exp$!approximation}%
  \index{approximation!of $\exp$}%
  If $\abs{x}\le 1$, then 
  \begin{equation}\label{eq:stima_exp}
    \abs{e^x - \sum_{k=0}^n \frac{x^k}{k!}}
    \le \sum_{k={n+1}}^{+\infty} \frac{\abs{x}^k}{k!}
    \le \frac{\abs{x}^{n+1}}{n\cdot n!}.
  \end{equation}
  In particular, setting $x=1$, we find
  \begin{equation}\label{eq:stima_e}
     0 < e - \sum_{k=0}^n \frac{1}{k!} \le \frac{1}{n \cdot n!}
  \end{equation}
  and for $n=5$, we obtain
  \begin{equation}
    2.716 < e < 2.719
  \end{equation}
  \end{theorem}
  %
  \begin{proof}
  The tail of the exponential series can be estimated 
  using the sum of a geometric series:
  \begin{align*}
    \abs{e^x - \sum_{k=0}^n \frac {x^k} {k!}}
    & \le \sum_{k=n+1}^{+\infty} \frac{\abs{x}^k}{k!} 
     = \sum_{k=0}^{+\infty}\frac{\abs{x}^{n+k+1}}{(n+k+1)!} \\
    & \le \sum_{k=0}^{+\infty}\frac{\abs{x}^{n+1+k}}{(n+1)^{k+1}\cdot n!}
          = \frac{\abs{x}^{n+1}}{(n+1)\cdot n!}\sum_{k=0}^{+\infty}
     \enclose{\frac{\abs{x}}{(n+1)}}^k \\
        & = \frac{\abs{x}^{n+1}}{(n+1)\cdot n!} \cdot
    \frac{1}{1-\frac{\abs{x}}{(n+1)}}
     = \frac{\abs{x}^{n+1}}{(n+1-\abs{x})\cdot n!}.
  \end{align*}
  If $\abs{x}\le 1$, we thus obtain the estimate~\eqref{eq:stima_exp}.
  If $x=1$, the exponential series has positive terms and thus 
  approximates $e$ from below: we thus obtain~\eqref{eq:stima_e}.
  For $n=5$, we have
  \[
   \sum_{k=0}^5 \frac{1}{k!} = 1 + 1 + \frac 1 2 + \frac{1}{6} + \frac {1}{24} + \frac{1}{120}
   = \frac{326}{120}
  \]
  Thus on one hand
  \[
    e \ge \frac{326}{120} \ge 2.716
  \]
  and on the other
  \[
   e \le \frac{326}{120} + \frac{1}{5\cdot 5!}
     \le 2.717 + 0.002 = 2.719
  \]
\end{proof}
  
\begin{theorem}[Irrationality of $e$]
\mymargin{$e\not\in \QQ$}%
\index{$e\not\in \QQ$}%
\mymark{**}%
\index{irrationality!of $e$}%
\index{$e$!is irrational}%
The number $e$ is irrational.
\end{theorem}
%
\begin{proof}
For every $n\in \NN$, we have 
\[
  n! \cdot e = \sum_{k=0}^{+\infty} \frac{n!}{k!}
   = \sum_{k=0}^n \frac{n!}{k!} + n!\sum_{k=n+1}^{+\infty} \frac{1}{k!}.
\]
The first term in the previous sum is an integer
since if $k\le n$, the ratio $n!/k!$ is an integer.
Thanks to theorem~\ref{th:approx_e}
for the second term, if $n>1$, we have:
\[
0 < n! \cdot \sum_{k=n+1}^{+\infty} \frac{1}{k!}
\le \frac{1}{n} < 1.
\]
It follows that for every $n\ge 2$, the number $n! e$
is strictly between two consecutive integers and thus is never 
an integer. Therefore, $e$ cannot be rational because 
if it were $e=\frac p q$, then $n! e \in \ZZ$ for every 
$n>q$.
\end{proof}